\chapter{Theoretical Background}
\label{ch:theory}

This chapter talks about the theoretical foundation for the modeling and control methods employed in this thesis.
It starts by discussing classical linear modeling, which serves as the baseline description of heat exchanger dynamics.
Subsequently, the Koopman operator framework is introduced as a data-driven approach for modeling nonlinear system behavior using lifted linear representations.

Model Predictive Control (MPC) is then introduced as the central control method implemented throughout this research.
The chapter also introduces the Kalman filter (KF), which is used for state estimation in both the linear-based and Koopman-based control schemes. 
Finally, the complete Koopman-based MPC structure is described, combining the Koopman prediction model, state estimation, and constrained optimal control into a single control framework.


\section{Koopman Operator}
\label{sec:Koopman_operator}

The Koopman operator framework provides a data-driven approach for modeling nonlinear dynamical systems using linear representations in an augmented state space.
Rather than modeling the nonlinear evolution of the system states directly, the Koopman approach transforms the measured variables into a higher-dimensional space, allowing the system dynamics to be approximated by a linear model.

Consider first an autonomous discrete-time nonlinear system described by
\begin{equation}
x_{k+1} = f(x_k),
\label{eq:autonomous_nonlinear_system}
\end{equation}
where $f$ represents a generally nonlinear function describing the system dynamics.

For controlled systems, the dynamics can be written as
\begin{equation}
x_{k+1} = f(x_k, u_k),
\label{eq:controlled_nonlinear_system}
\end{equation}
where $u_k$ denotes the control input. In this work, the control input is included later in the finite-dimensional lifted model through the input matrix $B_{\mathcal{K}}$.

A linear operator working on functions of the system state is known as the Koopman operator $\mathcal{K}$.
For a lifting function $g(\cdot)$, it satisfies
\begin{equation}
\mathcal{K} g(x_k) = g(x_{k+1}).
\end{equation}

A key characteristic of the Koopman operator is its linearity, despite the nonlinear nature of the underlying system dynamics. This enables the use of the operator's spectral properties to analyze system evolution \cite{mauroy2020koopman}.


The Koopman operator was originally proposed by Koopman \cite{koopman1931} and later developed further by Mezi'c \cite{mezic2005}. It is a linear operator that acts in an infinite-dimensional space. Since this is not practical for real applications, we use finite-dimensional approximations built from measured data. These approximations use a limited number of lifted states, chosen through a lifting function $g(\cdot)$ \cite{mauroy2020koopman}.

The lifting transformation $g(\cdot)$ maps the original system variables into a higher-dimensional space, such that the lifted state vector satisfies $z_k \in \mathbb{R}^{n_z}$, where typically $n_z \gg n_x$, with $n_x$ denoting the original state dimension.
As a result, the Koopman system matrices have dimensions
\begin{equation}
A_{\mathcal{K}} \in \mathbb{R}^{n_z \times n_z}, \quad
B_{\mathcal{K}} \in \mathbb{R}^{n_z \times n_u}, \quad
C_{\mathcal{K}} \in \mathbb{R}^{n_y \times n_z},
\end{equation}
where $n_u$ and $n_y$ denote the dimensions of the input and output, respectively.

As an example, a system with a single measured state ($n_x = 1$) can be expanded to a lifted state vector of dimension $n_z = 10$ by adding nonlinear combinations of the measured temperature and input signal. This results in a Koopman model with matrices $A_{\mathcal{K}} \in \mathbb{R}^{10 \times 10}$ and $B_{\mathcal{K}} \in \mathbb{R}^{10 \times 1}$, while preserving a linear structure suitable for prediction and control.
Selecting the lifted state dimension involves a trade-off between model accuracy and computational complexity. In practice, the dimension $n_z$ is determined by the chosen lifting function and identification method and is typically selected based on empirical performance and computational considerations \cite{korda2018}.

In the lifted space, the system dynamics can be approximated by a linear model suitable for prediction and control~\cite{korda2018} of the form
\begin{equation}
z_{k+1} = A_{\mathcal{K}} z_k + B_{\mathcal{K}} u_k,
\label{eq:koopman_state}
\end{equation}
\begin{equation}
\hat{y}_k = C_{\mathcal{K}} z_k.
\label{eq:koopman_output}
\end{equation}

A widely used technique for approximating the Koopman operator from data is Extended Dynamic Mode Decomposition (EDMD), which builds a finite-dimensional approximation based on a selected set of basis functions. The effectiveness of EDMD is highly sensitive to the choice of these basis functions \cite{williams2015}.

\section{Deep Koopman}
\label{sec:deep_koopman}

In this work, a Deep Koopman approach is used, where the lifting function is learned directly from data using a neural network~\cite{lusch2018deep,yeung2019learning,otto2019linearly}. This method enables the lifting to adapt to the measured system dynamics without requiring manual selection of basis functions.

The structure of the model is encoder–linear dynamics–decoder~\cite{goodfellow2016deep}. The measured output is mapped by the encoder to a lifted state
\begin{equation}
z_k = f_y(y_k),
\label{eq:koopman_encoder}
\end{equation}
where $f_y(\cdot)$ is a trainable nonlinear mapping. The encoder in the constructed model maps the measured output $y_k$ to a lifted state $z_k \in \mathbb{R}^{10}$ using a fully connected neural network with three hidden layers of 60 neurons each.


The evolution of the lifted state is described by the finite-dimensional Koopman model as in Equation~\eqref{eq:koopman_state}, where $A_{\mathcal K}$ represents the linear dynamics in the lifted space and $B_{\mathcal K}$ describes the influence of the control input.


The predicted output is obtained using a decoder,
\begin{equation}
\hat{y}_k = f_{\mathrm{dec}}(z_k).
\label{eq:koopman_decoder}
\end{equation}
In this work, the decoder is linear. Therefore, it corresponds directly to the output matrix $C_{\mathcal K}$ and the output equation can be written as in Equation~\eqref{eq:koopman_output}.

Figures~\ref{fig:encoder_network} and~\ref{fig:decoder_network} show schematic examples of the encoder and decoder used in this work.

\begin{figure}[!htbp]
    \centering
    \resizebox{\linewidth}{!}{%
    \begin{tikzpicture}[
        >=Latex,
        plain/.style={
            draw=none,
            fill=none,
            inner sep=0pt,
            minimum size=0pt
        },
        neuron/.style={
            draw,
            circle,
            minimum size=7mm,
            inner sep=0pt
        },
        dot/.style={
            draw,
            circle,
            fill=black,
            minimum size=1mm,
            inner sep=0pt
        },
        net/.style={
            matrix of nodes,
            nodes in empty cells,
            ampersand replacement=\&,
            column sep=12mm,
            row sep=2mm
        }
    ]
    
    \matrix[net] (mat) {
    |[plain]| \parbox{1.7cm}{\centering Input}
    \& |[plain]| \parbox{2.0cm}{\centering (60 neurons)}
    \& |[plain]| \parbox{2.0cm}{\centering (60 neurons)}
    \& |[plain]| \parbox{2.0cm}{\centering (60 neurons)}
    \& |[plain]| \parbox{2.0cm}{\centering (10 neurons)} \\
    
     \& |[neuron]| \& |[neuron]| \& |[neuron]| \& |[plain]|  \\
    |[plain]|  \& |[plain]|  \& |[plain]|  \& |[plain]|  \& |[plain]|  \\
    |[plain]|  \& |[neuron]| \& |[neuron]| \& |[neuron]| \& |[neuron]| \\
    |[plain]|  \& |[plain]|  \& |[plain]|  \& |[plain]|  \& |[plain]|  \\
    |[plain]|  \& |[dot]|    \& |[dot]|    \& |[dot]|    \& |[neuron]| \\
    |[plain]|  \& |[dot]|    \& |[dot]|    \& |[dot]|    \& |[dot]|    \\
    |[neuron]| \& |[neuron]| \& |[neuron]| \& |[neuron]| \& |[dot]|    \\
    |[plain]|  \& |[plain]|  \& |[plain]|  \& |[plain]|  \& |[plain]|  \\
    |[plain]|  \& |[neuron]| \& |[neuron]| \& |[neuron]| \& |[neuron]| \\
    |[plain]|  \& |[plain]|  \& |[plain]|  \& |[plain]|  \& |[plain]|  \\
    |[plain]|  \& |[neuron]| \& |[neuron]| \& |[neuron]| \& |[plain]|  \\
    };
    
    % input arrow
    \draw[->] ([xshift=-12mm]mat-8-1.west) -- (mat-8-1.west)
    node[midway, above] {$y_k$};
    
    % input -> hidden layer 1
    \foreach \r in {2,4,8,10,12}{
        \draw[->, gray!65] (mat-8-1) -- (mat-\r-2);
    }
    
    % hidden layer 1 -> hidden layer 2
    \foreach \r in {2,4,8,10,12}{
        \foreach \s in {2,4,8,10,12}{
            \draw[->, gray!45] (mat-\r-2) -- (mat-\s-3);
        }
    }
    
    % hidden layer 2 -> hidden layer 3
    \foreach \r in {2,4,8,10,12}{
        \foreach \s in {2,4,8,10,12}{
            \draw[->, gray!45] (mat-\r-3) -- (mat-\s-4);
        }
    }
    
    % hidden layer 3 -> latent layer
    \foreach \r in {2,4,8,10,12}{
        \foreach \s in {4,6,10}{
            \draw[->, gray!45] (mat-\r-4) -- (mat-\s-5);
        }
    }
    
    % labels on right of latent layer
    \node[right=6mm of mat-4-5] {$z_1$};
    \node[right=6mm of mat-6-5] {$z_2$};
    \node[right=10mm of mat-7-5] {$\vdots$};
    \node[right=6mm of mat-10-5] {$z_{10}$};
    
    \end{tikzpicture}%
    }
    \caption{Schematic example of the encoder architecture used in this work. The measured output $y_k$ is mapped by a fully connected neural network with three hidden layers of 60 neurons each into a 10-dimensional latent state $z_k$.}
    \label{fig:encoder_network}
\end{figure}


\begin{figure}[!htbp]
    \centering
    \begin{tikzpicture}[
        >=Latex,
        plain/.style={
            draw=none,
            fill=none,
            inner sep=0pt,
            minimum size=0pt
        },
        neuron/.style={
            draw,
            circle,
            minimum size=7mm,
            inner sep=0pt
        },
        dot/.style={
            draw,
            circle,
            fill=black,
            minimum size=1mm,
            inner sep=0pt
        },
        net/.style={
            matrix of nodes,
            nodes in empty cells,
            column sep=18mm,
            row sep=2mm
        }
    ]
    
    \matrix[net] (mat) {
    |[plain]| \parbox{2.0cm}{\centering (10 neurons)} &
    |[plain]| \parbox{1.8cm}{\centering Output} \\
    
     & |[plain]|  \\
    |[plain]|  & |[plain]|  \\
    |[neuron]| & |[plain]|  \\
    |[plain]|  & |[plain]|  \\
    |[dot]|    & |[plain]|  \\
    |[dot]|    & |[plain]|  \\
    |[neuron]| & |[neuron]| \\
    |[plain]|  & |[plain]|  \\
    |[neuron]| & |[plain]|  \\
    };
    
    % connections latent -> output
    \foreach \r in {4,8,10}{
        \draw[->, gray!45] (mat-\r-1) -- (mat-8-2);
    }
    
    % labels on left of latent layer
    \node[left=3mm of mat-4-1] {$z_1$};
    \node[left=6mm of mat-6-1] {$z_2$};
    \node[left=3mm of mat-8-1] {$\vdots$};
    \node[left=2mm of mat-10-1] {$z_{10}$};
    
    % output arrow
    \draw[->] (mat-8-2.east) -- ([xshift=12mm]mat-8-2.east)
    node[midway, above] {$\hat{y}_k$};
    
    \end{tikzpicture}
    \caption{Schematic example of the decoder architecture used in this work. The 10-dimensional lifted state is mapped linearly to the predicted outlet temperature $\hat{y}_k$.}
    \label{fig:decoder_network}
\end{figure}

\FloatBarrier

All parts of the model, including the encoder, the lifted linear dynamics, and the decoder, are identified simultaneously from measured trajectories. The identification is performed using the Neuromancer framework, which formulates the training as a structured optimization problem and enables gradient-based learning.

After training, the learned matrices $A_{\mathcal K}$, $B_{\mathcal K}$, and $C_{\mathcal K}$ define a linear state-space model in the lifted coordinates. This model can then be used directly as the prediction model in MPC.

The individual steps of the identification process are summarized below.

\begin{enumerate}
    \item Experimental data are collected from the heat exchanger by applying step changes to the input and measuring the outlet temperature.

    \item The measured signals are preprocessed and scaled, and organized into input--output sequences.

    \item A nonlinear lifting function is learned using a neural network encoder, which maps the measured variables into a lifted state space according to Equation~\eqref{eq:koopman_encoder}.

    \item In the lifted space, the system dynamics are described by the linear controlled Koopman model given in Equation~\eqref{eq:koopman_state}.


    \item The lifted state is mapped back to the output space using the decoder in Equation~\eqref{eq:koopman_decoder}. In the implemented model, this decoder is linear and corresponds to the output matrix in Equation~\eqref{eq:koopman_output}.


    \item All model components are identified jointly using multi-step prediction over measured trajectories, ensuring that the predicted output in Equation~\eqref{eq:koopman_output} matches the measured outlet temperature.


\end{enumerate}

To ensure accurate representation of system dynamics across the prediction horizon, the model is trained with a combination of loss functions applied to the predicted trajectories.

The main loss term is the multi-step output prediction error
\begin{equation}
\mathcal{L}_{y} =
\sum_{k=1}^{N_{\mathrm{p}}}
\left\| y_k - \hat{y}_k \right\|_2^2,
\end{equation}
which penalizes deviations between measured and predicted outputs over the full prediction horizon $N_{\mathrm{p}}$ used during training.

In addition, a one-step prediction loss is included
\begin{equation}
\mathcal{L}_{\text{1-step}} =
\left\| y_1 - \hat{y}_1 \right\|_2^2,
\end{equation}
to improve local prediction accuracy.

To ensure consistency between the lifted linear dynamics and the encoded states, a latent-space loss is introduced
\begin{equation}
    \mathcal{L}_{x} =
    \sum_{k=1}^{N_{\mathrm{p}}}
    \left\| z_k^{\mathrm{enc}} - z_k^{\mathrm{pred}} \right\|_2^2,
    \end{equation}
where $z_k^{\text{enc}}$ denotes the encoded lifted state and $z_k^{\text{pred}}$ is obtained from the linear Koopman dynamics.

A reconstruction loss is also included to ensure that the encoder--decoder pair preserves the original measurements
\begin{equation}
\mathcal{L}_{\text{rec}} =
\left\| y_0 - \hat{y}_0 \right\|_2^2.
\end{equation}

The overall optimization problem is then formulated as
\begin{equation}
\min_{\theta, A_{\mathcal{K}}, B_{\mathcal{K}}, C_{\mathcal{K}}}
\quad
\alpha_y \mathcal{L}_{y}
+ \alpha_{\mathrm{1step}} \mathcal{L}_{\mathrm{1step}}
+ \alpha_x \mathcal{L}_{x}
+ \alpha_{\mathrm{rec}} \mathcal{L}_{\mathrm{rec}},
\label{eq:koopman_training_loss}
\end{equation}
where $\theta$ denotes the trainable parameters of the nonlinear encoder. The coefficients
$\alpha_y$, $\alpha_{\mathrm{1step}}$, $\alpha_x$, and $\alpha_{\mathrm{rec}}$ are weighting parameters used to balance the individual loss terms.

The model parameters are optimized using gradient-based methods over measured trajectories. 
This multi-objective formulation ensures accurate long-term prediction, consistency of the lifted dynamics, and faithful reconstruction of the measured outputs.

Once trained, the neural network is used only to define the lifting transformation. 
The resulting Koopman model is then represented as a linear state-space system in the lifted coordinates and used for prediction and control.

\newpage
\section{Kalman Filter}
\label{sec:KF}

The Kalman filter (KF) is a state-estimation method that uses a mathematical model and noisy measurements to estimate the internal states of dynamic systems. In control applications, it provides reliable state estimates for feedback controllers, particularly when not all system states are directly measurable \cite{kalman1960,morari1999}. In Koopman-based MPC, the lifted states are usually not directly measurable and therefore require an estimate for feedback control~\cite{valabek2026experimental}.

Consider a discrete-time linear system with process and measurement noise described by
\begin{equation}
x_{k+1} = A x_k + B u_k + w_k,
\label{eq:strejc_state}
\end{equation}
\begin{equation}
y_k = C x_k + v_k,
\label{eq:strejc_state2}
\end{equation}
where $w_k$ and $v_k$ denote zero-mean Gaussian process and measurement noise, respectively.


The KF estimates the system state by minimizing the mean-square error, assuming the presence of Gaussian noise. The algorithm operates in two main phases: prediction and correction. Here, $\hat{x}_{k|k}$ refers to the state estimate at time $k$ based on all measurements up to that point, while $\hat{x}_{k|k-1}$ is the predicted state estimate before the current measurement is taken into account.

During the prediction phase, the state estimate and its associated covariance are propagated using the system model,
\begin{equation}
\hat{x}_{k|k-1} = A \hat{x}_{k-1|k-1} + B u_{k-1},
\end{equation}
\begin{equation}
P_{k|k-1} = A P_{k-1|k-1} A^\top + Q_{\mathrm{KF}},
\end{equation}
where $P_{k|k-1}$ represents the prediction error covariance matrix, $Q_{\mathrm{KF}}$ is the process noise covariance matrix of the Kalman filter, and $u_{k-1}$ denotes the control input applied at the previous sampling instant.

In the correction phase, the predicted state estimate is updated using the latest measurement,
\begin{equation}
K_k = P_{k|k-1} C^\top \left( C P_{k|k-1} C^\top + R_{\mathrm{KF}} \right)^{-1},
\end{equation}
\begin{equation}
\hat{x}_{k|k} = \hat{x}_{k|k-1} + K_k \left( y_k - C \hat{x}_{k|k-1} \right),
\end{equation}
\begin{equation}
P_{k|k} = (I - K_k C) P_{k|k-1},
\end{equation}
where $P_{k|k}$ is the updated estimation error covariance matrix, $R_{\mathrm{KF}}$ is the measurement noise covariance matrix of the KF, $K_k$ is the Kalman gain, and $I \in \mathbb{R}^{n \times n}$ is the identity matrix.

The state estimations required by the MPC controller in our work are obtained by the Kalman filter. To estimate the Koopman state vector utilized for prediction and control in the Koopman-based MPC, the KF operates in the lifted state space. In contrast, for the linear MPC, the KF estimates the internal state of the first-order linear model that describes the main thermal dynamics. A fair and consistent foundation for comparison is ensured by using the same estimating structure for both approaches.


\section{Model Predictive Control}
\label{sec:MPC}

Model Predictive Control (MPC) is an advanced control strategy based on repeatedly solving a constrained finite-horizon optimization problem \cite{rawlings_mpc,martin2019mpc}. At each sampling instant, a prediction model is used to predict future system behavior over a predefined prediction horizon, and an optimal sequence of control inputs is computed.

MPC follows the receding horizon principle, where an optimal sequence of control inputs is calculated for the prediction horizon, but only the first input is actually implemented on the system \cite{martin2019mpc}. At each sampling step, the optimization problem is resolved with the latest state estimates.

In this work, MPC is formulated for discrete-time linear systems described by the state-space model
\begin{equation}
x_{k+1} = A x_k + B u_k,
\end{equation}
\begin{equation}
y_k = C x_k + D u_k,
\end{equation}
where $x_k \in \mathbb{R}^n$ denotes the system state vector, $u_k \in \mathbb{R}^m$ is the control input, and $y_k \in \mathbb{R}^p$ is the system output. The matrices $A$, $B$, $C$, and $D$ describe the discrete-time prediction model. In many thermal systems, including the system considered in this work, the direct feedthrough term $D$ is zero or negligible.

The simple linear MPC was obtained from a first-order step-response model. The continuous-time transfer function can be written as
\begin{equation}
G(s) = \frac{K}{\tau s + 1},
\label{eq:first_order_transfer}
\end{equation}
where $K$ is the steady-state gain and $\tau$ is the time constant. Using zero-order-hold discretization with sampling period $\Ts$, the corresponding discrete-time state-space model is obtained as
\begin{equation}
A = e^{-\Ts/\tau}, \qquad
B = K\left(1 - e^{-\Ts/\tau}\right), \qquad
C = 1, \qquad
D = 0.
\label{eq:first_order_discrete_matrices}
\end{equation}
These matrices define the prediction model used by the MPC controller based on the linear step-response approximation.

The control objective is to drive the system to a desired steady-state operating point. This is formulated using a quadratic cost function that penalizes both deviations of the predicted output from zero and excessive control input. At each sampling instant, MPC solves the following optimization problem:


    
    \begin{equation}
        \begin{aligned}
        \min_{\{u_k\}_{k=0}^{N-1}} \quad &
        \sum_{k=0}^{N-1}
        \left(
        \|y_k\|_Q^2 + \|u_k\|_R^2
        \right) \\
        \text{s.t.} \quad
        & x_{k+1} = A x_k + B u_k, \\
        & y_k = C x_k, \\
        & u_{\min} \le u_k \le u_{\max}, \\
        & y_{\min} \le y_k \le y_{\max}, \\
        & k = 0,\ldots,N-1, \\
        & x_{0} = \hat{x}_{t},
        \end{aligned}
        \label{eq:mpc_optimization}
    \end{equation}
    

    where $N$ denotes the prediction horizon, and $\hat{x}_{t}$ is the estimated system state at the current sampling instant $t$. The index $k = 0,\ldots,N-1$ denotes the prediction step within the horizon. The weighting matrices $Q \succeq 0$ and $R \succ 0$ determine the relative importance of output regulation and control effort, respectively. $Q$ increases the penalty on output deviations, while $R$ penalizes large or aggressive control moves.


 
% In this thesis, $N = 20$ was selected based on empirical results.

Input and output constraints capture the physical and operational boundaries of the system. For thermal processes, input constraints usually correspond to actuator saturation, while output constraints ensure that temperatures remain within safe operating ranges \cite{rawlings_mpc, martin2019mpc}.
The choice of prediction horizon $N$ influences both the control performance and the computational load. In practice, $N$ is selected to achieve a compromise between prediction accuracy and the real-time requirements of the application \cite{rawlings_mpc}.

This MPC formulation is used consistently for both the linear-based and Koopman-based prediction models. The difference between the two control approaches lies solely in the prediction model matrices used in the optimization problem, enabling a direct and fair comparison.





\section{Koopman-Based Model Predictive Control}

Koopman-based Model Predictive Control (Koopman MPC) integrates the conventional first-order MPC approach with a data-driven Koopman model and uses a Kalman filter for state estimation. Unlike linear MPC, which relies on a first-order model, Koopman MPC formulates its prediction model in the lifted state space described in Section~\ref{sec:deep_koopman}. The lifted-state dynamics and output mapping from equations~\eqref{eq:koopman_state} and~\eqref{eq:koopman_output} are incorporated directly into the MPC prediction process \cite{korda2018,koopman_mpc_review}.

From the MPC perspective, the Koopman-based prediction model is treated as a standard linear state-space system. The same quadratic cost function and input/output constraints used in first-order MPC are applied. The key distinction is that the prediction model operates on the lifted state vector $z_k$, and a state estimator is used to provide the initial state for prediction.

At each sampling instant $t$, the Koopman MPC controller solves the following finite-horizon optimization problem
\begin{equation}
\begin{aligned}
\min_{\{u_k\}_{k=0}^{N-1}} \quad &
\sum_{k=0}^{N-1}
\left(
\|y_k\|_Q^2 + \|u_k\|_R^2
\right) \\
\text{s.t.} \quad
& z_{k+1} = A_{\mathcal{K}} z_k + B_{\mathcal{K}} u_k, \\
& y_k = C_{\mathcal{K}} z_k, \\
& u_{\min} \le u_k \le u_{\max}, \\
& y_{\min} \le y_k \le y_{\max}, \\
& k = 0,\ldots,N-1, \\
& z_0 = \hat{z}_t,
\end{aligned}
\label{eq:koopman_mpc_optimization}
\end{equation}

where $\hat{z}_t \in \mathbb{R}^{n_z}$ denotes the estimated lifted state vector at the current sampling instant $t$, obtained from the Kalman filter and used as the initial condition for the MPC prediction. The index $k = 0,\ldots,N-1$ denotes the prediction step within the horizon. The output constraints are applied to the predicted physical output, which is computed from the lifted state using the output mapping described in Equation~\eqref{eq:koopman_output}.

In closed-loop operation, the Koopman MPC controller proceeds as follows.
At each sampling instant, the heat exchanger outlet temperature is measured and scaled.
This measurement, along with the most recent control input, is used by the Kalman filter described in Section~\ref{sec:KF} to estimate the lifted state vector $\hat{z}_k$.
The estimated state $\hat{z}_k$ is then used as the initial condition of the MPC prediction model.

The optimization problem is solved over the finite prediction horizon, but only the first control input from the computed sequence is implemented on the system. At each sampling instant, this process is repeated with the latest measurements and updated state estimates.


\begin{figure}[htbp]
    \centering
    \def\svgwidth{0.75\textwidth}
    \input{svg-inkscape/SecDIag.drawio_svg-tex.pdf_tex}
    \caption{Simplified block diagram of the Koopman-based MPC control structure.}
    \label{fig:koopman_mpc_structure}
  \end{figure}
  

Figure~\ref{fig:koopman_mpc_structure} illustrates the closed-loop structure, where the Kalman filter provides an estimate of the lifted state $\hat{z}_k$ from measured outputs, and the MPC controller computes the control input based on this estimate.

